\subsection{乘法自同态}

设 X 是局部紧拓扑空间，\(\mu\) 是 X 上的正测度。设 \(p \in [1, +\infty[\) 。

设 \(g \in \mathcal{L}^{\infty}(\mathrm{X},\mu)\) 。用 \(m_{g}\) 表示由 \(f \mapsto g \cdot f\) 所给出的从 \(\mathcal{L}^{p}(\mathrm{X},\mu)\) 到自身的映射。映射 \(m_{g}\) 是线性的且连续的，因为

\[
\mathrm{N} _ {p} (m _ {g} (f)) \leqslant \mathrm{N} _ {\infty} (g) \mathrm{N} _ {p} (f)\tag{1}
\]

对每个函数 \(f \in \mathcal{L}^{p}(\mathrm{X}, \mu)\) 成立。特别地，函数 \(m_{g}(f)\) 是 \(\mu\)-零测的，如果 f 是 \(\mu\)-零测的。于是，映射 \(m_{g}\) 通过过渡到商诱导出 \(\mathrm{L}^{p}(\mathrm{X}, \mu)\) 的一个自同态，我们将它记作 \(\widetilde{m}_{g}\) 。

引理 7. --- 设 \(g \in \mathcal{L}^{\infty}(\mathrm{X},\mu)\) 。自同态 \(\tilde{m}_g\) （\(\mathrm{L}^{p}(\mathrm{X},\mu)\) 的）是单射的，当且仅当由 \(x \in \mathrm{X}\) （满足 \(g(x) = 0\) ）所成的集是局部 \(\mu\)-零测的。

用 A 表示由 \(x \in X\) （满足 \(g(x) = 0\) ）所成的集。

假设 A 是局部 \(\mu\)-零测的。设 \(f \in \mathcal{L}^p(\mathrm{X}, \mu)\) ，\(\widetilde{f}\) 是它在 \(\mathrm{L}^p(\mathrm{X}, \mu)\) 中的类。假设 \(\widetilde{m}_g(\widetilde{f}) = 0\) 。按定义，这等价于函数 \(fg\) 是 \(\mu\)-零测的，于是由 \(x \in \mathrm{X}\) （满足 \(f(x)g(x) \neq 0\) ）所成的集 B 是 \(\mu\)-零测的。函数 \(f\) 在 \(\mathrm{A} \cup \mathrm{B}\) 之外为零，故是 \(\mu\)-零测的。这就蕴含自同态 \(\widetilde{m}_g\) 是单射的。

反之，假设 A 不是局部 \(\mu\)-零测的。设 C 是 X 的紧子集，使得 A∩C 不是 \(\mu\)-零测的，并设 \(\varphi\) 是 A∩C 的特征函数。函数 \(\varphi\) 在 L \(^{p}\) (X, \(\mu\) ) 中的类不为零，但 \(m_{g}(\varphi)\) 的类为零，故 \(m_{g}\) 不是单射的。

由于局部 \(\mu\)-零测函数与 \(\mu\)-零测函数之积是 \(\mu\)-零测的，自同态 \(\widetilde{m}_g\) 只依赖于 \(g\) 在 \(\mathrm{L}^{\infty}(\mathrm{X},\mu)\) 中的类。对 \(\widetilde{g}\in \mathrm{L}^{\infty}(\mathrm{X},\mu)\) ，我们也用 \(\widetilde{m}_{\widetilde{g}}\) 记自同态 \(\widetilde{m}_g\) （\(\mathrm{L}^p (\mathrm{X},\mu)\) 的），这里 \(g\in \mathcal{L}^{\infty}(\mathrm{X},\mu)\) 是以 \(\widetilde{g}\) 为类的任一函数。我们称 \(\widetilde{m}_{\widetilde{g}}\) 为由 \(\widetilde{g}\) 所作的、\(\mathrm{L}^p (\mathrm{X},\mu)\) 中的乘法自同态。

命题 5. --- 设 \(p \in [1, +\infty[\) 。映射 \(m: g \mapsto \widetilde{m}_g\) （从 \(\mathrm{L}^{\infty}(\mathrm{X}, \mu)\) 到 \(\mathcal{L}(\mathrm{L}^{p}(\mathrm{X}, \mu))\) 中的）是含幺巴拿赫代数的一个等距态射。

对 \(a\) ， \(b\) （属于 \(\mathbf{C}\) ）以及 \(g_{1}\) ， \(g_{2}\) （属于 \(\mathrm{L}^{\infty}(\mathrm{X},\mu)\) ），可直接验证 \(\widetilde{m}_{ag_1+bg_2}=a\widetilde{m}_{g_1}+b\widetilde{m}_{g_2}\) ，故映射 \(\widetilde{m}\) 是线性的。此外，我们有 \(\widetilde{m}_{1}=1_{\mathrm{L}^{p}(\mathrm{X},\mu)}\) 与 \(\widetilde{m}_{g_{1}g_{2}}=\widetilde{m}_{g_{1}}\widetilde{m}_{g_{2}}\) ，故映射 \(\widetilde{m}\) 是从 \(\mathrm{L}^{\infty}(\mathrm{X},\mu)\) 到 \(\mathcal{L}(\mathrm{L}^{p}(\mathrm{X},\mu))\) 中的含幺代数态射。

上面的公式 (1) 表明 \(\widetilde{m}\) 的范数 \(\leqslant 1\) 。

设 \(g \in \mathcal{L}^{\infty}(\mathrm{X},\mu)\) ，\(\widetilde{g}\) 是它在 \(\mathrm{L}^{\infty}(\mathrm{X},\mu)\) 中的类。设 \(\varepsilon > 0\) 。由 \(x \in X\) （满足 \(|g(x)| > \| \widetilde{g}\|_{\infty} - \varepsilon\) ）所成的集 Y 不是局部 \(\mu\)-零测的。因此存在 X 的紧子集 C 使得 \(\mu (\mathrm{Y} \cap \mathrm{C}) > 0\) 。设 \(\varphi\) 是 \(\mathrm{Y} \cap \mathrm{C}\) 的特征函数。由于

\[
\begin{array}{c} \| \widetilde {m} _ {\widetilde {g}} (\varphi) \| _ {p} = \Bigl (\int_ {\mathrm{X}} | \varphi g | ^ {p}   d \mu \Bigr) ^ {1 / p} \geqslant \Bigl (\int_ {\mathrm{Y}} | \varphi g | ^ {p}   d \mu \Bigr) ^ {1 / p} \\ \geqslant (\| \widetilde {g} \| _ {\infty} - \varepsilon) \mu (\mathrm{Y} \cap \mathrm{C}) ^ {1 / p} = (\| \widetilde {g} \| _ {\infty} - \varepsilon) \| \varphi \| _ {p}, \end{array}
\]

便得到 \(\|\widetilde{m}_{\widetilde{g}}\|\geqslant\|\widetilde{g}\|_{\infty}-\varepsilon\) 。由于 \(\varepsilon\) 是任意的，可推出 \(\|\widetilde{m}_{\widetilde{g}}\|\geqslant\|\widetilde{g}\|_{\infty}\) ，这就证明了 \(\widetilde{m}\) 是等距的。

引理 8. --- 设 \((g_n)\) 是 \(\mathcal{L}^\infty(\mathrm{X},\mu)\) 中 \(\mu\)-几乎处处逐点收敛的有界序列。设 \(\widetilde{m} \in \mathrm{L}^\infty(\mathrm{X},\mu)\) 是其极限的类。那么 \(\widetilde{m}_{g_n}\) 收敛于 \(\widetilde{m}_{\widetilde{g}}\) ，这里是在赋有逐点收敛拓扑的空间 \(\mathcal{L}(\mathrm{L}^p(\mathrm{X},\mu))\) 中。

设 \(f \in \mathcal{L}^p(\mathrm{X},\mu)\) 。序列 \((g_{n}f)\) 在 \(\mathcal{L}^p(\mathrm{X},\mu)\) 中有界，且 \(\mu\)-几乎处处逐点收敛于 \(gf\) 。由勒贝格定理（INT, IV, p. 137, § 3, no. 7，定理 6），序列 \((g_{n}f)\) 收敛于 \(gf\) （在 \(\mathcal{L}^p(\mathrm{X},\mu)\) 中），断言得证。

现在考虑 p = 2 的情形。

命题 6. --- 映射 m: \(g \mapsto \widetilde{m}_{g}\) 是从星代数 \(\mathrm{L}^{\infty}(\mathrm{X}, \mu)\) 到星代数 \(\mathcal{L}(\mathrm{L}^{2}(\mathrm{X}, \mu))\) 中的含幺等距态射。

特别地，对每个 \(g \in \mathrm{L}^{\infty}(\mathrm{X}, \mu)\) ，乘法自同态 \(\widetilde{m}_{g}\) 是正规的；它是埃尔米特的，当且仅当 g \(\mu\)-几乎处处局部取实值。

由命题 5，映射 \(\tilde{m}\) 是从 \(\mathrm{L}^{\infty}(\mathrm{X},\mu)\) 到 \(\mathcal{L}(\mathrm{L}^{2}(\mathrm{X},\mu))\) 中的含幺巴拿赫代数的单射等距态射。设 \(g\in \mathcal{L}^{\infty}(\mathrm{X},\mu)\) 。对 \(f_{1}\) 与 \(f_{2}\in \mathcal{L}^{2}(\mathrm{X},\mu)\) ，我们有

\[
\langle f _ {1} \mid \widetilde {m} _ {g} (f _ {2}) \rangle = \int_ {\mathrm{X}} \overline {{f _ {1} (x)}} g (x) f _ {2} (x) d \mu (x) = \langle \widetilde {m} _ {\overline {{g}}} (f _ {1}) \mid f _ {2} \rangle ,
\]

由此得出 \(\widetilde{m}_g^* = \widetilde{m}_{\overline{g}}\) ，这就证明了 \(m\) 是对合态射。最后的断言由此得出（参见 I, p. 106，命题 5）。

推论. --- 设 \(g \in \mathrm{L}^{\infty}(\mathrm{X},\mu)\) 。

\begin{enumerate}
\def\labelenumi{\alph{enumi})}
\item
  \(\widetilde{m}_g\) 在 \(\mathcal{L}(\mathrm{L}^2 (\mathrm{X},\mu))\) 中的谱是 \(\mu\)-本质意义下 \(g\) 的像；
\item
  自同态 \(\widetilde{m}_g\) 是正的，当且仅当 \(g\) \(\mu\)-几乎处处局部为正；
\item
  对任意函数 \(f \in \mathcal{C}(\mathrm{Sp}(\widetilde{m}_g))\) ，我们有 \(f(\widetilde{m}_g) = \widetilde{m}_{f \circ g}\) 。
\end{enumerate}

由于 \(\widetilde{m}\) 是单射的，由 I, p. 106 的命题 5，我们有 \(\mathrm{Sp}(\widetilde{m}_{g}) = \mathrm{Sp}(g)\) ；结论由 IV, p. 185 的引理 6、IV, p. 185 的命题 4 以及 I, p. 115 的定义 6 得出。

命题 7. --- 态射 \(\tilde{m}\) （从 \(\mathrm{L}^{\infty}(\mathrm{X},\mu)\) 到 \(\mathcal{L}(\mathrm{L}^{2}(\mathrm{X},\mu))\) 中的）的像是代数 \(\mathcal{L}(\mathrm{L}^{2}(\mathrm{X},\mu))\) 的一个极大交换子代数。

只须证明：若 \(u \in \mathcal{L}(\mathrm{L}^{2}(\mathrm{X}, \mu))\) 是与 \(\widetilde{m}_{g}\) （对每个函数 \(g \in \mathcal{L}^{\infty}(\mathrm{X}, \mu)\) ）都交换的自同态，则 u 属于 \(\widetilde{m}\) 的像。设 u 是这样一个自同态。

我们用 \(\tilde{f}\) 表示 \(\mathrm{L}^2 (\mathrm{X},\mu)\) 中函数 \(f\in \mathcal{L}^2 (\mathrm{X},\mu)\) 的类。用 \(v\) 表示从 \(\mathcal{L}^2 (\mathrm{X},\mu)\) 到 \(\mathrm{L}^2 (\mathrm{X},\mu)\) 中的、由 \(f\mapsto u(\widetilde{f})\) 所定义的线性映射。我们有 \(v\circ m_g = \widetilde{m}_g\circ v\) （对所有 \(g\in \mathcal{L}^{\infty}(\mathrm{X},\mu)\) ）。

对 X 的任意 \(\mu\)-可测子集 Y，我们用 \(\varphi_{\mathrm{Y}}\) 表示它的特征函数；自同态 \(\widetilde{m}_{\varphi_{\mathrm{Y}}}\) 是 \(\mathrm{L}^2 (\mathrm{X},\mu)\) 的一个正交投影算子。

首先假设 X 是紧的，从而常值函数 1 的类属于 \(\mathrm{L}^{2}(\mathrm{X},\mu)\) 。设 g 是 \(\mathcal{L}^{2}(\mathrm{X},\mu)\) 中的函数，它在 \(\mathrm{L}^{2}(\mathrm{X},\mu)\) 中的类为 \(v(1)\) 。

设 \(c\) 是正实数，Y 是由 \(x \in \mathrm{X}\) （满足 \(|g(x)| \geqslant c\) ）所成的集；它是一个 \(\mu\)-可测集。我们在 \(\mathrm{L}^2(\mathrm{X}, \mu)\) 中有等式 \(\widetilde{\varphi_{\mathrm{Y}}} g = \widetilde{m}_{\varphi_{\mathrm{Y}}}(v(1)) = v(m_{\varphi_{\mathrm{Y}}}(1)) = v(\varphi_{\mathrm{Y}})\) 。

又由于 \(|c\varphi_{\mathrm{Y}}| \leqslant |g\varphi_{\mathrm{Y}}|\) ，我们得到

\[
c ^ {2} \mu (\mathrm{Y}) = \int_ {\mathrm{X}} (c \varphi_ {\mathrm{Y}}) ^ {2} d \mu \leqslant \int_ {\mathrm{X}} | \varphi_ {\mathrm{Y}} g | ^ {2} d \mu = \| \widetilde {u} (\varphi_ {\mathrm{Y}}) \| ^ {2} \leqslant \| u \| ^ {2} \mu (\mathrm{Y}).
\]

这个不等式表明 \(\mu(\mathrm{Y}) = 0\) （若 \(c > \|u\|\) ），从而函数 \(g\) \(\mu\)-几乎处处以 \(\|u\|\) 为界。最后，对每个函数 \(f \in \mathcal{C}(\mathrm{X})\) ，我们有

\[
u (\widetilde {f}) = v (m _ {f} (1)) = \widetilde {m} _ {f} (v (1)) = \widetilde {m} _ {f} (\widetilde {g}) = \widetilde {m} _ {g} (\widetilde {f}),
\]

由此得到等式 \(u = \tilde{m}_{g}\) ，特别地 \(\mathrm{N}_{\infty}(g) = \| u\|\) 。

现在考虑一般情形。存在 X 的两两不相交紧子集所成的局部可数族 \((\mathrm{X}_{i})_{i\in\mathrm{I}}\) ，使得 \(Z = X - \bigcup_{i \in I} X_i\) 是局部 \(\mu\)-零测的（INT, IV, p.190, § 5, no. 9，命题 14）。设 \(\mu_i\) 是 \(\mu\) 在 \(X_i\) 上诱导的测度（INT, IV, p. 186, § 5, no. 7，定义 4）。

由 IV, p. 179 的命题 1，对每个 \(i \in I\) ，像 \(\mathrm{E}_i\) （\(\widetilde{m}_{\varphi_{\mathrm{X}_i}}\) 的）与 \(\mathrm{L}^2(\mathrm{X}_i, \mu_i)\) 等同——通过 \(f \mapsto f|\mathrm{X}_i\) 。对每个函数 \(g \in \mathcal{L}^\infty(\mathrm{X}, \mu)\) ，乘法自同态 \(\widetilde{m}_g\) 与 \(\widetilde{m}_{\varphi_{\mathrm{X}_i}}\) 交换，因而使 \(\mathrm{E}_i\) 保持稳定，且 \(\mathrm{E}_i\) 的、由 \(\widetilde{m}_g\) 过渡到子空间而得的自同态与乘以 \(g|\mathrm{X}_i\) 的乘法自同态（\(\mathrm{L}^2(\mathrm{X}_i, \mu_i)\) 上的）重合。

由于 \(u\) 与 \(m_{\varphi_{X_i}}\) 交换，它也使子空间 \(\mathrm{E}_i\) 保持稳定。用 \(u_i\) 表示 \(\mathrm{L}^2 (\mathrm{X}_i,\mu_i)\) 的、由 \(u\) 过渡到子空间而得的自同态。

由于从 \(\mathcal{L}^{\infty}(\mathrm{X},\mu)\) 到 \(\mathcal{L}^{\infty}(\mathrm{X}_{i},\mu_{i})\) 中的、由 \(g\mapsto g|\mathrm{X}_{i}\) 所定义的映射是满射的，假设蕴含 \(u_{i}\) 与 \(\widetilde{m}_{g}\) 交换（对每个函数 \(g\in \mathcal{L}^{\infty}(\mathrm{X}_{i},\mu_{i})\) ）。由证明的第一部分，存在函数 \(g_{i}\in \mathcal{L}^{\infty}(\mathrm{X}_{i},\mu_{i})\) ，使得 \(u_{i} = \widetilde{m}_{g_{i}}\) 且 \(\mathrm{N}_{\infty}(g_{i})\leqslant \| u_{i}\| \leqslant \| u\|\) 。

X 上与 \(g_{i}\) 在 \(X_{i}\) 上重合、在 Z 上为零的函数 g 是有界且 \(\mu\)-可测的（INT, IV, p. 193, § 5, no. 10，命题 16），故 \(g \in \mathcal{L}^{\infty}(X, \mu)\) 。对每个 \(i \in I\) ，u 在 \(E_{i}\) 上的限制与 \(\tilde{m}_{g}\) 的限制重合；因此，由 IV, p. 179 的命题 1, c)，我们有 \(u = m_{g}\) 。

推论. --- 设 u 是希尔伯特空间 \(\mathrm{L}^{2}(\mathrm{X},\mu)\) 的自同态，它与 \(\widetilde{m}_{g}\) （对每个函数 \(g\in\mathcal{K}(\mathrm{X})\) ）交换。那么存在唯一的元素 \(f\in\mathrm{L}^{\infty}(\mathrm{X},\mu)\) ，使得 \(u=\widetilde{m}_{f}\) 。

由命题 7，只须证明 \(u\) 与 \(\widetilde{m}_{g}\) 交换（对所有 \(g\in\mathcal{L}^{\infty}(\mathrm{X},\mu)\) ）。设 \(h_{1}\) 与 \(h_{2}\) 是 \(\mathcal{L}^{2}(\mathrm{X},\mu)\) 中的元素，它们的类 \(\widetilde{h}_{1}\) 与 \(\widetilde{h}_{2}\) 属于 \(\mathrm{L}^{2}(\mathrm{X},\mu)\) 。设 \(k_{1}\) （相应地， \(k_{2}\) ）是 \(\mathcal{L}^{2}(\mathrm{X},\mu)\) 中的函数，其类为 \(u(\widetilde{h}_{1})\) （相应地， \(u^{*}(\widetilde{h}_{2})\) ）。

设 \(h = h_1 \overline{k}_2 - k_1 \overline{h}_2\) ；我们有 \(h \in \mathcal{L}^1(\mathrm{X}, \mu)\) 。定义测度 \(\nu = h \cdot \mu\) （在 \(\mathrm{X}\) 上）；它是有界的。对每个 \(g \in \mathcal{L}^\infty(\mathrm{X}, \mu)\) ，我们有

\[
\begin{array}{c} \langle \widetilde {h} _ {2} | u (\widetilde {m} _ {g} (\widetilde {h} _ {1})) - \widetilde {m} _ {g} (u (\widetilde {h} _ {1})) \rangle = \langle u ^ {*} (\widetilde {h} _ {2}) | \widetilde {m} _ {g} (\widetilde {h} _ {1}) \rangle - \langle \widetilde {h} _ {2} | \widetilde {m} _ {g} (u (\widetilde {h} _ {1})) \rangle \\ = \int_ {\mathrm{X}} g \cdot h _ {1} \cdot \overline {{k}} _ {2} d \mu - \int_ {\mathrm{X}} g \cdot k _ {1} \cdot \overline {{h}} _ {2} d \mu = \nu (g). \end{array}
\]

于是，按假设有 \(\nu(g) = 0\) （对每个函数 \(g \in \mathcal{K}(\mathrm{X})\) ），即 \(\nu = 0\) 。由此可得

\[
\langle \widetilde {h} _ {2} | u (\widetilde {m} _ {g} (\widetilde {h} _ {1})) - \widetilde {m} _ {g} (u (\widetilde {h} _ {1})) \rangle = \nu (g) = 0
\]

对所有 \(g \in \mathcal{L}^{\infty}(\mathrm{X},\mu)\) 成立。由于这对所有 \(h_1\) 与 \(h_2\) （\(\mathcal{L}^2 (\mathrm{X},\mu)\) 中的）都成立，我们有 \(u\circ \widetilde{m}_g = \widetilde{m}_g\circ u\) 。

注. --- 在下文中，我们常径直用 \(m_{g}\) 表示乘以 g 的乘法算子（\(\mathrm{L}^{p}(\mathrm{X},\mu)\) 上的）。

